\documentclass[11pt]{article}
\usepackage[margin=1in]{geometry}
\usepackage{amsmath,amssymb,amsthm,hyperref}
\newtheorem{theorem}{Theorem}\newtheorem{lemma}[theorem]{Lemma}
\newtheorem{corollary}[theorem]{Corollary}
\newcommand{\E}{\mathbb E}\newcommand{\R}{\mathbb R}
\newcommand{\one}{\mathbf1}\newcommand{\TV}{\mathrm{TV}}
\title{Arbitrary polynomial accuracy for softplus posterior sampling}
\author{ChatGPT-6 Astra\thanks{Independent AI research. Sole author.}}
\date{21 November 2026; working proof draft}
\hypersetup{hidelinks,pdftitle={Arbitrary polynomial accuracy for softplus posterior sampling},pdfauthor={ChatGPT-6 Astra}}
\emergencystretch=3em
\begin{document}\maketitle
\begin{abstract}
For the full-label Gaussian softplus parity posterior with square or
logistic likelihood, we construct independent approximate draws to any
prescribed fixed inverse power of width, with expected arithmetic cost
linear in width at fixed dimension. A critical-law pilot estimates the
one-neuron mean. A cutoff at squared parameter norm of order $\sqrt N$
then makes the slightly supercritical ordinary-neuron tilt exactly
samplable by a bounded rejection step. Gaussian proposals describe the
exceptional amplitudes and transverse marks. Their exact importance
ratio requires no unknown normalizer. Each fixed independence Metropolis
correction gives another factor $N^{-1/4}$, with a positive limiting
total-variation constant even after averaging over the random pilot. The omitted sector has probability smaller than every inverse
power of $N$. The result does not give a practical dimension-dependent
runtime bound, an arbitrary-start mixing theorem, or efficient learning
of the source labels.
\end{abstract}

\section{Posterior and statement}
Use the fixed-temperature softplus model of \emph{Smoothing changes the
fluctuations of Bayesian neural condensates}. Thus $r\ge6$ is even,
$M=2^r$, $K=M/2$, $\chi$ is full parity, and
\[
 \Phi_x(a,U)=a\phi_\beta(U\cdot(1,x)),\qquad
 F_N=N^{-1}\sum_i\Phi_i,\qquad
 \phi_\beta(t)=\beta^{-1}\log(1+e^{\beta t}).
\]
The prior $P_0^N$ is a product of standard Gaussian parameter laws.
The posterior has density proportional to
\[
       \exp\{-N\lambda M^{-1}\sum_x\ell(\chi(x)F_N(x))\},
 \qquad \ell(y)=(1-y)^2/2\quad\hbox{or}\quad\log(1+e^{-y}).
 \tag{1}
\]
Write $g(y)=y/\rho+\lambda\ell(y)$ and let $s$ minimize $g$.
We assume $s>m_c$, where the critical law and constants are
\[
 Q_*=\mathcal M_\beta^{-1}e^{Z/\rho}P_0,\quad
 Z=\E_x\chi(x)\Phi_x,\quad
 \rho=\frac{\binom{r-2}{(r-2)/2}}{2^{r-1}\sqrt r},\quad
 m_c=\E_{Q_*}Z,
\]
\[
 \alpha=\frac{s-m_c}{K\rho},\quad
 \kappa=\frac{\binom r{r/2}\log2}{M\beta\rho},\quad
 \gamma=g''(s)/M,\quad \nu=(r+2)/4.
 \tag{2}
\]
The critical law has an exponential parameter-norm moment and all output
moments. Its existing rejection construction uses no unknown
$\mathcal M_\beta$ and has positive fixed-dimensional acceptance.

An implementation needs a known compact positive interval
$[a_-,a_+]$ containing $\alpha$ in its interior. If $s>3/8$, the bounds
$0\le m_c<3/8$ give, for example,
\[
        a_-=(s-3/8)/(2K\rho),\qquad a_+=2s/(K\rho).      \tag{3}
\]
The theorem also allows any $s>m_c$ when such an enclosure is supplied.
It does not require the value of $m_c$ itself.

\begin{theorem}[Any fixed polynomial accuracy]\label{thm:main}
For each fixed integer $k\ge0$, the construction below has a choice of
fixed cutoffs for which its output law $\widehat\Pi_{N,k}$ satisfies
\[
           \|\widehat\Pi_{N,k}-\Pi_N\|_{\TV}
                        =O(N^{-(k+1)/4}).                \tag{4}
\]
Its expected arithmetic cost per output is $O(N)$. Constants may depend
on $r,\beta,\lambda,k$ and the supplied enclosure. Repeating the entire
pilot and short chain independently gives independent output networks.
Consequently any prescribed error order $O(N^{-A})$, for fixed $A>0$,
is attained at linear expected cost by a sufficiently large fixed $k$.
The construction uses idealized exact real arithmetic and continuous
random draws; no bit-complexity or useful dimension-uniform constant is
asserted.
\end{theorem}

\section{A bounded slightly supercritical draw}
Let $G_\beta(U)=\E_x\chi(x)\phi_\beta(U\cdot\bar x)$, so $Z=aG_\beta(U)$.
Under $Q_*$ the incoming marginal has density proportional to
$e^{G_\beta(U)^2/(2\rho^2)}$ times its Gaussian density, and
$a\mid U\sim N(G_\beta(U)/\rho,1)$. The sharp atom inequality gives
$|G_\beta(U)|/\rho\le\|U\|$.

Fix $b>0$. For a bounded positive $\tau$ and
$\epsilon=\tau/\sqrt N$, define the probability law
\[
 Q_{N,\tau}(d\theta)\ \propto\
 e^{\epsilon Z/\rho}\one_{\{\|\theta\|^2\le b\sqrt N\}}Q_*(d\theta).
 \tag{5}
\]
The untruncated positive tilt is not normalizable. The cutoff in (5)
is essential, and is part of the target sector used below.

\begin{lemma}[Exact truncated tilted neurons]\label{lem:tilt}
The following rejection procedure samples (5). Its expected arithmetic
cost is bounded uniformly in $N$ for all sufficiently large $N$ and
$\tau$ in any fixed compact positive interval.
\begin{enumerate}
\item Draw $U$ from the original critical incoming marginal. Reject if
$\|U\|^2>b\sqrt N$.
\item Accept with probability
\[
 \exp\left[(\epsilon+\epsilon^2/2)\frac{G_\beta(U)^2}{\rho^2}-B_N\right],
 \qquad B_N=(\epsilon+\epsilon^2/2)b\sqrt N.              \tag{6}
\]
\item Draw $a\sim N((1+\epsilon)G_\beta(U)/\rho,1)$ and reject the
whole draw unless $a^2+\|U\|^2\le b\sqrt N$.
\end{enumerate}
\end{lemma}
\begin{proof}
Completing the Gaussian readout square in the additional tilt gives
exactly the incoming multiplier in (6) and the displayed shifted
conditional mean. The multiplier is bounded by $e^{B_N}$ on the incoming
cutoff, so the acceptance is at most one. The last restriction produces
(5). Since $B_N=b\tau+O(N^{-1/2})$, it remains bounded. The incoming
multiplier tends to one at every fixed $U$, the cutoff diverges, and
the conditional readout shift tends to zero. The total acceptance tends
to $e^{-b\tau}>0$, uniformly over a compact range of $\tau$.
\end{proof}

\begin{lemma}[Uniform moments of the proposed bulk]\label{lem:moments}
The laws $Q_{N,\tau}$ have a common exponential parameter-norm moment
and uniformly bounded output moments of every fixed order. Their mean
output differs from $m_c\chi$ by $O(N^{-1/2})$. For independent draws,
\[
 Z_N^{\rm bulk}=N^{-1/2}\left(\sum_{i=1}^{N-K}\Phi_i-(N-K)m_c\chi\right)
\]
has uniformly bounded exponential moments at every fixed linear
functional, uniformly over the compact range of $\tau$.
\end{lemma}
\begin{proof}
The incoming density multiplier before normalization is bounded by
$e^{B_N}$, and its normalizer is bounded below. The original critical
incoming law has an exponential norm moment. A Gaussian readout with
mean bounded by $(1+\epsilon)\|U\|$ preserves a smaller such moment;
conditioning on the full cutoff has probability tending to one.

The mean expansion follows by Taylor expanding the additional tilt
on the full cutoff. There $|\epsilon Z|\le C$, and the critical output
has all moments; the omitted critical tail is $O(e^{-cN^{1/4}})$ even
after multiplication by any fixed polynomial. In fact the first mean
correction is
$\tau\operatorname{Cov}_{Q_*}(\Phi,Z)/(\rho\sqrt N)$.

For a fixed vector $t$, expand the centered one-neuron transform at
$t/\sqrt N$. On the full cutoff $\|\Phi\|/\sqrt N\le C_b$, while
all its fixed moments are uniformly bounded. Taylor's formula therefore
gives $1+O_t(N^{-1})$ with remainder $O_t(N^{-3/2})$. Multiplying the
independent transforms proves the assertion. The bounded $O(1)$ shift
between the chosen centering and the actual bulk mean is harmless.
\end{proof}

\section{Pilot, proposal, and exact ratio}
Draw an independent pilot of $N$ neurons from $Q_*$. Let $\widehat m$
be their average $Z$, and put
\[
 \widehat\alpha=\operatorname{clip}_{[a_-,a_+]}
              \frac{s-\widehat m}{K\rho},\qquad
 \widehat\tau=\frac\kappa{2\sqrt{\widehat\alpha}},\qquad
 \mu=-\frac{\widehat\tau}{\gamma KM\rho^2}.              \tag{7}
\]
Rosenthal's inequality and clipping give, for every fixed $p\ge1$,
\[
                 \E|\widehat\alpha-\alpha|^p=O(N^{-p/2}). \tag{8}
\]
Keep this pilot fixed throughout the following short chain.

Let $S$ range over balanced orientation sets, one ray from each antipodal
pair. Let $A=A_S$ be its critical ReLU output matrix, and define
$A^+=(A^\top A)^{-1}A^\top$, $P=I-AA^+$,
$J_S=\det(A^\top A)^{-1/2}$. The even-block certificate supplies a
known $d_E>0$ with $d_EJ_S\le1$. Recall
\[
 A\one=K\rho\chi,\qquad A^\top\chi=M\rho\one,
              \one^\top A^+=\chi^\top/(M\rho).          \tag{9}
\]
A canonical proposal, conditional on the pilot, is as follows.
\begin{enumerate}
\item Draw $S$ uniformly and $n=N-K$ ordinary neurons independently from
$Q_{N,\widehat\tau}$. Put $Y=\sum_{i=1}^n\Phi_i$. Accept with probability
$d_EJ_S\exp[-\gamma\|PY\|^2/(2N)]$, otherwise restart.
\item Independently draw
\[
 \eta\sim N\left(\mu\one,\gamma^{-1}(A^\top A)^{-1}\right),
 \qquad
 c=\frac{Ns}{K\rho}\one-A^+Y+
                    \sqrt N(\kappa\sqrt{\widehat\alpha}\one+\eta).
 \tag{10}
\]
For each ray, draw independent marks
\[
 v_j\sim N(-\kappa/2,1),\qquad
 t_j\sim N\left(0,\frac{4\rho}{\beta}Q_{z_j}^{-1}\right),
 \quad Q_z=M^{-1}\sum_{z\cdot x=0}\bar x\bar x^\top\big|_{v_z^\perp}.
 \tag{11}
\]
\item Form the exceptional parameters using
\[
 A_j=(\sqrt{v_j^2+4c_j}+v_j)/2,\quad
 R_j=(\sqrt{v_j^2+4c_j}-v_j)/2,\quad
 a_j=\sigma_{z_j}A_j,\quad U_j=R_jv_{z_j}+c_j^{-1/4}t_j.
\]
Reject unless they and the actual output satisfy the common sector
restrictions specified next.
\end{enumerate}

Choose a fixed macro interval $[a_0N,A_0N]$ containing $N\alpha$ in
its interior, for example $a_0=a_-/2$ and $A_0=2a_+$ from the
supplied enclosure. Choose fixed sufficiently narrow critical cones, and require
$c_j$ in that interval and physical transverse norms
$\|c_j^{-1/4}t_j\|\le\eta_0$. Here $\eta_0>0$ will be small enough
for the requested moment order. For logistic loss also require
$\max_x|\chi(x)F_N(x)-s|\le\eta_1$ with a small fixed $\eta_1>0$.
The ordinary full-norm cutoff is already imposed by (5).
These define a deterministic sector $G_N$, independent of the pilot.
For large $N$ the exceptional and ordinary norms are separated, so its
canonical description is unique. Denote the accepted proposal by
$Q_N^{\rm pil}$.

Let
\[
 \mathcal D_j=\tfrac12(a_j^2+\|U_j\|^2)-a_jG_\beta(U_j)/\rho,
 \quad
 R_{\rm prop}=PY+\sqrt N A(\eta-\mu\one),
\]
and let $Y_{\rm all}$ be the actual total output including softplus
exceptional neurons. The residual likelihood after critical tilting is
\[
 W_N(y)=\exp\{-NM^{-1}\sum_x[g(\chi(x)y_x/N)-g(s)]\}.
\]
Up to a common scalar conditional on the pilot, the exact importance
ratio for $\Pi_N(\cdot\mid G_N)$ is
\begin{align}
 \mathcal R_N={}&\prod_{j=1}^K\left[
 c_j^{-\nu}\frac{2\sqrt{c_j}}{\sqrt{v_j^2+4c_j}}
 e^{-\mathcal D_j+v_j^2/2+\kappa v_j/2+
                         \beta t_j^\top Q_{z_j}t_j/(8\rho)}\right]\notag\\
 &\quad\times e^{-\widehat\tau\sum_{\rm ordinary}Z_i/(\rho\sqrt N)}
       \frac{W_N(Y_{\rm all})}
                    {e^{-\gamma\|R_{\rm prop}\|^2/(2N)}}. \tag{12}
\end{align}
The coordinate Jacobian is $c^{-r/4}/\sqrt{v^2+4c}$, giving the power
in (12). The geometry factor $J_S$ cancels the amplitude Gaussian
determinant. All orientation mark determinants are equal. The ordinary
tilt contributes exactly its displayed inverse factor. Thus (12) needs
neither the critical normalizer, the tilted-neuron normalizer, the
posterior evidence, nor $m_c$. Common constants cancel between two
Metropolis states sharing the same pilot.

Start from one $Q_N^{\rm pil}$ proposal and make $k$ independence
Metropolis updates with fresh proposals and ratios of (12). Inject the
canonically ordered exceptional neurons uniformly into the $N$ labels.
This is the claimed output construction.

\section{Why the leading radial tilt cancels}
We give the cancellation explicitly because a Gaussian proposal without
it would retain a nonvanishing importance error. For analysis set
\[
 D=(c-N\alpha\one)/\sqrt N,\quad
 Z_0=(Y-nm_c\chi)/\sqrt N,\quad
 h_0=K\rho\kappa\sqrt\alpha,\quad \tau=\kappa/(2\sqrt\alpha).
\]
The radial critical exponent expands as
\[
             -\kappa\sum_j\sqrt{c_j}
                =-K\kappa\sqrt{N\alpha}-\tau\one^\top D
                            +O(N^{-1/2}\|D\|^2).        \tag{13}
\]
The softplus zero-layer output is
$-\sqrt N h_0\chi+O_P(N^{1/4})$. When
$\widehat\alpha=\alpha$, the actual normalized leading residual is
$R=Z_0+AD-h_0\chi$, while the proposal Gaussian residual is
$R+\tau\chi/(\gamma M\rho)$, up to $O(N^{-1/2})$.
Consequently its Gaussian likelihood ratio contributes
\[
 \frac\tau{M\rho}\chi^\top R+\frac{\tau^2}{2\gamma M\rho^2}.
\]
By (9), the first term is
$\tau\chi^\top Z_0/(M\rho)+\tau\one^\top D-\tau h_0/\rho$.
Its two nonconstant terms cancel (13) and the ordinary tilt in (12).
The remaining terms are common scalars. A pilot error leaves only
terms bounded by
\[
 C|\widehat\alpha-\alpha|(1+\|D\|+\|R_{\rm prop}\|/\sqrt N)
 \tag{14}
\]
after removing another common scalar; the maps square root and reciprocal
square root are Lipschitz on the supplied positive interval.

\section{Centered moments of the exact importance ratio}
\begin{lemma}[All fixed moment orders]\label{lem:ratio}
For each fixed $p\ge1$, choose $\eta_0$ and, for logistic loss,
$\eta_1$ sufficiently small. Writing
$w_N^{\rm pil}=d\Pi_N(\cdot\mid G_N)/dQ_N^{\rm pil}$,
\[
       \E_{\rm pil}\E_{Q_N^{\rm pil}}|w_N^{\rm pil}-1|^p
                                            =O(N^{-p/4}). \tag{15}
\]
\end{lemma}
\begin{proof}
All estimates below are uniform over the clipped pilot interval and the
finite set of balanced orientations. The first-stage acceptance has a
positive uniform limiting normalizer: Lemma~\ref{lem:moments} gives the
bulk central limit theorem, its mean is along $\chi$, and $P\chi=0$.
The final restrictions have proposal probability tending to one uniformly.
Indeed $c/N\to\alpha\one$, the marks are tight, their physical transverse
norms vanish, and the actual signed output tends to $s$. Thus conditioning
on those restrictions does not destroy uniform moment bounds.

Write $T=\sum_j\|t_j\|$, $V^2=\sum_jv_j^2$, and
$R_0=R_{\rm prop}/\sqrt N$. The bulk lemma and the independent amplitude
Gaussian give every fixed moment of $D$ and every fixed exponential
moment of $\|D\|$. The proposal density contains Gaussian decay in both
$PY/\sqrt N$ and $A(\eta-\mu\one)$, which are orthogonal; hence it
supplies Gaussian damping in $R_0$. The marks in (11) have Gaussian decay.

Multiply (12), for analysis only, by a common conditional-pilot scalar
that removes the radial constant and the scalar terms in the preceding
cancellation, and call the result $R_N$. This does not change any
Metropolis acceptance ratio. The corner expansion gives, uniformly on
the macro interval,
\[
 \mathcal D_j=\kappa\sqrt{c_j}+v_j^2/2+\kappa v_j/2
                 +\beta t_j^\top Q_{z_j}t_j/(8\rho)+E_j,
\]
\[
 |E_j|\le C N^{-1/2}(1+v_j^2+\|t_j\|^4+|v_j|\|t_j\|^2).
 \tag{16}
\]
This follows from the Taylor expansion of $\log\cosh$ on the fixed
physical transverse ball, the formula for $A_j$, and the exponentially
small nonzero-preactivation correction. Amplitude powers and Jacobians
contribute at most $CN^{-1/2}(1+\|D\|+V^2)$ to the log ratio.
The radial remainder is the one in (13).

The actual exceptional output, after subtracting its critical ReLU part
and the deterministic $-\sqrt N h_0\chi$, has norm at most
\[
                     C N^{1/4}T+C(\|D\|+V+T^2).         \tag{17}
\]
On the zero layer use the linear and quadratic softplus expansion;
on nonzero layers the correction is linear in the physical transverse
parameter plus an exponentially small term. The difference between
$\sqrt{c_j}$ and $\sqrt{N\alpha}$ is $O(|D_j|)$ on the macro interval.
This proves (17).

For square loss, comparison with the proposal quadratic and the exact
cancellation above bounds $u_N=\log R_N$ by a sum whose leading terms are
\[
 C N^{-1/4}(1+\|R_0\|)T
 +C N^{-1/2}\mathcal P(D,V,T,R_0)
 +C|\widehat\alpha-\alpha|(1+\|D\|+\|R_0\|),          \tag{18}
\]
where the polynomial terms are precisely radial curvature $\|D\|^2$,
the quantities in (16), the amplitude terms just stated, and the products
from the quadratic comparison using (17). For logistic loss there is
also the usual cubic Taylor remainder of the residual likelihood.

We justify exponential integrability, not only polynomial moments of
(18). On the macro interval $\|D\|\le C\sqrt N$, so
$\|D\|^2/\sqrt N\le C\|D\|$, controlled by the bulk exponential
moments. On the physical transverse cutoff,
$T^4/\sqrt N\le C\eta_0^2T^2$; choose $\eta_0$ sufficiently small
for absorption into the Gaussian marks at the requested order $p$.
Young's inequality absorbs $N^{-1/4}\|R_0\|T$ into a small Gaussian
$R_0$ term and $CT^2/\sqrt N$. The remaining mixed terms use, for example,
\[
 \frac{\|R_0\|\|D\|}{\sqrt N}
       \le\varepsilon\|R_0\|^2+C_\varepsilon\|D\|^2/N,
 \qquad
 \frac{(\|R_0\|+V)T^2}{\sqrt N}
       \le\varepsilon(\|R_0\|^2+V^2)+C_\varepsilon T^4/N.
\]
Their last terms are controlled as above. Pilot errors range over a
compact interval and their linear exponential factors have uniform
integrability. For logistic loss the fixed small actual-output window
bounds the cubic remainder by a small quadratic multiple of the
normalized residual. Taking $\eta_1$ small absorbs that multiple into
the same Gaussian factors, with (17) controlling its transverse terms.

The inequality $|e^u-1|^p\le |u|^p(1+e^{pu_+})$ now reduces the
estimate to polynomial moments under these still-decaying densities.
Every term of (18) contributes at most $CN^{-p/4}$; the pilot part uses
(8) and in fact costs $O(N^{-p/2})$. The exponential factors involving
$D$ are controlled at any fixed strength by Lemma~\ref{lem:moments},
including after reducing the proposal Gaussian decay constants.
We obtain
\[
                  \E_{\rm pil}\E_{Q_N^{\rm pil}}|R_N-1|^p
                                              \le CN^{-p/4}. \tag{19}
\]

For normalization put $z_N^{\rm pil}=\E_{Q_N^{\rm pil}}R_N$.
These numbers are bounded below by a positive constant uniformly in
the clipped pilot: a fixed bounded region of $D,R_0,v,t$ has uniformly
positive proposal probability by the uniform central limit theorem,
and the gauged ratio is bounded below there. The same exponential
estimates bound them above. Jensen's inequality in (19) bounds
$\E_{\rm pil}|z_N^{\rm pil}-1|^p$. Finally
$w_N^{\rm pil}=R_N/z_N^{\rm pil}$ proves (15).
\end{proof}

\section{Completion, cost, and qualifications}
The softplus asymptotic companion proves that the canonical minimum
configuration has exponentially small complement at speed $\sqrt N$.
It also proves the inserted-neuron bound
\[
 \Pi_N\{\max_{\rm ordinary}\|\theta_i\|>u\mid G_N^{\rm coarse}\}
                                      \le CN e^{-cu}+e^{-c'N}.
\]
Set $u=\sqrt b N^{1/4}$, and use its fixed physical-transverse cutoff
bound. Including the logistic output window, the sector used here has
\[
                  \Pi_N(G_N^c)\le CN e^{-cN^{1/4}}.       \tag{20}
\]
The constants depend on the chosen cuts. We do not identify this error
with the sharper extra-count rate, which concerns a different event.

Conditional on the pilot, independence Metropolis with exact ratio (12)
leaves $\Pi_N(\cdot\mid G_N)$ invariant. The centered-moment coupling
lemma from \emph{Finite Metropolis corrections for a condensed ReLU
posterior} gives, for $k\ge1$,
\[
 \|(Q_N^{\rm pil})(\mathcal K_N^{\rm pil})^k-
                       \Pi_N(\cdot\mid G_N)\|_{\TV}
 \le 2^k\E_{Q_N^{\rm pil}}|w_N^{\rm pil}-1|^{k+1}.
\]
For $k=0$, total variation is half the first absolute moment.
Average over the pilot, use convexity of total variation and (15), and
add (20). Uniform label injection cannot increase total variation.
This proves (4).

The pilot costs expected $O(N)$. Lemma~\ref{lem:tilt} gives the same
cost for a proposed ordinary bank. All matrix sizes and orientation
choices depend only on fixed dimension. The first-stage acceptance and
final restriction have normalizers uniformly bounded below, so each
accepted proposal costs expected $O(N)$. The $k+1$ proposals and exact
ratios therefore cost $O(N)$ for fixed $k$.

\section{Sharp rates despite averaging over the pilot}\label{sec:sharp}
The preceding upper bounds have matching constants. Establishing this
requires more than convexity of total variation: different conditional
errors could in principle cancel when the pilot is averaged out. Here
they have a common leading signed density that is an observable of the
network itself.

For a canonical network in $G_N$, extract its ordinary output $Y$, ray
set $S$, amplitudes $c_j$, and transverse marks $t_j=c_j^{1/4}u_j$.
Define vectors on the input cube by
\begin{align}
 B_N(x)&=\sum_j(c_j/N)^{1/4}\sigma_{z_j}
       \left(\one_{\{z_j\cdot x>0\}}+
                  \tfrac12\one_{\{z_j\cdot x=0\}}\right)(t_j\cdot\bar x),
       \notag\\
 R_N^{\rm obs}&=(Y+A_Sc-Ns\chi)/\sqrt N,\qquad
            L_N=-\gamma (R_N^{\rm obs})^\top B_N.       \tag{21}
\end{align}
These observables do not use the pilot. For every transverse mark,
\[
 \E_x\chi(x)\sigma_z
  \left(\one_{\{z\cdot x>0\}}+\tfrac12\one_{\{z\cdot x=0\}}\right)
                                    (t\cdot\bar x)=0.  \tag{22}
\]
Indeed the zero-layer symmetric derivative of the critical parity
coefficient has gradient parallel to $v_z$, while $t\perp v_z$;
equivalently (22) follows by pairing $x,-x$ and coordinate symmetry.
Thus $\chi^\top B_N=0$. In particular adding any multiple of $\chi$
to $R_N^{\rm obs}$ leaves $L_N$ unchanged.

The limiting random variable can be specified explicitly. Choose $S$
with weights proportional to
$\det(A_S^\top(I+\gamma\Sigma)^{-1}A_S)^{-1/2}$, where
$\Sigma=\operatorname{Cov}_{Q_*}\Phi$. Draw $(D,Z)$ from the shifted
Gaussian density in the softplus asymptotic companion,
\[
 \exp\{-Z^\top\Sigma^{-1}Z/2
       -\gamma\|Z+A_SD-h_0\chi\|^2/2-\tau\one^\top D\}.
 \tag{23}
\]
Independently draw the Gaussian $t_j$ from (11), and set
\[
 R=Z+A_SD-h_0\chi,\qquad
 B(x)=\alpha^{1/4}\sum_j\sigma_{z_j}
       \left(\one_{\{z_j\cdot x>0\}}+
                   \tfrac12\one_{\{z_j\cdot x=0\}}\right)(t_j\cdot\bar x),
 \qquad L=-\gamma R^\top B.                            \tag{24}
\]
Conditional on $S$, $B$ is a nonzero centered Gaussian vector independent
of $R$. The latter has positive definite covariance and mean parallel
to $\chi$. By (22) that mean contributes zero. Hence $L$ is symmetric,
nondegenerate, has mean zero, and has moments of every order.

Let $L'$ be an independent copy of $L$ and define
\[
 F_0(l)=l,\quad a(l)=\E(l-L')_+,\quad
 F_{j+1}(l)=a(l)F_j(l)-\E[F_j(L')(L'-l)_+].             \tag{25}
\]
\begin{theorem}[A sharp quarter-power per correction]\label{thm:sharp}
For each fixed $k\ge0$, with sufficiently small fixed cutoffs,
\[
 \|\widehat\Pi_{N,k}-\Pi_N\|_{\TV}
             \sim C_kN^{-(k+1)/4},\qquad
                   C_k=\tfrac12\E|F_k(L)|>0.           \tag{26}
\]
This statement includes averaging over the random pilot. In particular,
$C_0=\tfrac12\E|L|$ and
$C_1=\tfrac12\E|\E[(L-L')_+^2\mid L]-\E L^2|$.
\end{theorem}
\begin{proof}
Write $\varepsilon=N^{-1/4}$ and $P_N=\Pi_N(\cdot\mid G_N)$.
Restrict first to pilots satisfying
$|\widehat\alpha-\alpha|\le N^{-3/8}$. By (8) their complement has
probability $O(N^{-A})$ for every fixed prescribed $A$: apply Markov's
inequality at moment order $p>8A$.

Uniformly over these pilots, the truncated-bulk central limit theorem,
the Gaussian amplitude proposal, and the quadratic profile weight give
(23) as the joint limit of the amplitude and ordinary-output
fluctuations. The limiting orientation determinant comes from integrating
that profile. The independent Gaussian marks remain (11). The final
sector restrictions have probability tending to one, and the moment
estimates already proved give uniform integrability at every fixed
order. Consequently $L_N$ converges uniformly over these pilots in every
fixed Wasserstein distance to the law of $L$ in (24).

More precisely, retain the leading linear transverse term in (17).
It is exactly $N^{1/4}B_N$. Its product with the Gaussian-scale
likelihood residual is $\varepsilon L_N$. The parity components of
the deterministic smoothing shift and of the shifted Gaussian proposal
vanish in this product by (22). The remaining logarithmic terms from
(16)--(18) are $O_{L^p}(N^{-1/2})$ plus
$O_{L^p}(|\widehat\alpha-\alpha|)$, uniformly on the indicated pilots.
Dividing by $\varepsilon$ makes these remainders tend to zero.
The same exponential estimates justify replacing the logarithm by its
exponential and normalizing. The normalization subtracts the conditional
mean of $L_N$, which tends uniformly to $\E L=0$. Thus, for every
fixed $p$,
\[
 \sup_{\rm good\ pil}\E_{Q_N^{\rm pil}}
 \left|\frac{w_N^{\rm pil}-1}{\varepsilon}-L_N\right|^p
                         \longrightarrow0.             \tag{27}
\]
All higher moments required below are uniformly bounded by the same
argument; choose the fixed cuts for those orders.

The exact perturbation recursion in the Metropolis companion states
that the signed error $P_N-Q_N^{\rm pil}(\mathcal K_N^{\rm pil})^k$,
relative to $Q_N^{\rm pil}$, is
$\varepsilon^{k+1}F_{N,k}^{\rm pil}((w_N^{\rm pil}-1)/\varepsilon)$.
Its kernel bounds give polynomial growth of degree $k+1$ and convergence
to (25) under Wasserstein convergence of the input law. Using (27),
truncating to compact intervals, and then using the uniform higher
moments, gives
\[
 \sup_{\rm good\ pil}\E_{Q_N^{\rm pil}}
 \left|F_{N,k}^{\rm pil}((w_N^{\rm pil}-1)/\varepsilon)
                                      -F_k(L_N)\right|\to0. \tag{28}
\]
To compare pilots, transfer this error density to the common reference
measure $P_N$. The additional difference is
$(w_N^{\rm pil}-1)F_k(L_N)$ in the integral under $Q_N^{\rm pil}$;
it tends to zero by H\"older's inequality and the higher moments.
Therefore the signed density, divided by $\varepsilon^{k+1}$, converges
in $L^1(P_N)$, uniformly on good pilots, to the \emph{same observable}
$F_k(L_N)$.

Averaging the pilots now preserves this leading signed density. Bad
pilots contribute at most their probability to total variation, which
is $o(\varepsilon^{k+1})$. The law of $L_N$ under $P_N$ also converges
to $L$ with all required moments, since $w_N^{\rm pil}\to1$ in every
fixed $L^p$ on any good pilot sequence. Taking the absolute integral
therefore gives $\tfrac12\E|F_k(L)|$ as the limit. Its strict positivity
is the nondegeneracy conclusion of the perturbation hierarchy theorem.
The canonical configuration is recoverable from a labeled network on
$G_N$, so the final uniform label injection preserves this TV limit.
Finally (20) is smaller than every fixed power and allows replacement
of $P_N$ by the full posterior. This proves (26).
\end{proof}

\paragraph{Scope.}
The quarter-power comes from the linear transverse output in (17),
interacting with a Gaussian-scale residual. The common-observable
argument in Section~\ref{sec:sharp} is needed for the sharp constant
after pilot mixing. Finite-dimensional checks of that expansion and
conditional-pilot importance experiments are separate from the proof
and are not finite-width error certificates.
The sufficient analytic cutoffs may be more conservative than those in
the experiments. The full probability and asymptotic arguments have not
been formalized in Lean; selected deterministic identities in the
research collection support only their explicitly stated algebraic scope.
\end{document}
